\documentclass{article}
\usepackage{amsmath,amssymb,amsthm}
\newtheorem{theorem}{Theorem}
\begin{document}
\title{Three heads transfer every two-head strict sign function}
\author{}\date{}\maketitle

\paragraph{Verification status and scope.}
The theorem below is a complete written finite-domain proof, with an
exact rational checker. It is not Lean-checked. It gives a constant-size
orientation transfer and does not establish an asymptotic saving.
It preserves the input dimension and uses the existing admissibility
criterion for positive affine denominators with strictly uniform slopes.
It neither asserts an analogous constant-factor transfer for arbitrarily
many heads nor supplies exact axis cancellation.

\begin{theorem}
Let $X$ be a finite subset of a Boolean cube, and let $A_1,A_2,B_1,B_2$
be affine forms with $B_1,B_2>0$ on the entire cube. Suppose
$F=A_1/B_1+A_2/B_2$ has no zero on $X$.
Choose any affine form $L$ positive on the entire cube whose slopes
are all strictly positive or all strictly negative.
For every sufficiently small $\varepsilon>0$, the three-head score
\[
 G_\varepsilon=
 \frac{A_1/\varepsilon}{L-\varepsilon B_2}
 +\frac{A_2/\varepsilon}{L-\varepsilon B_1}
 -\frac{(A_1+A_2)/\varepsilon}{L}
 \tag{1}
\]
has the same strict signs as $F$ on $X$. All three denominators are
positive on the entire cube, have the orientation of $L$, and converge
coefficientwise to $L$.
\end{theorem}
\begin{proof}
Put $P=A_1B_2+A_2B_1$. Exact common-denominator arithmetic gives
\[
 G_\varepsilon=
 \frac{LP-\varepsilon(A_1+A_2)B_1B_2}
 {L(L-\varepsilon B_1)(L-\varepsilon B_2)}.
 \tag{2}
\]
In particular $G_\varepsilon\to P/L^2$ at every cube vertex.
The signs of $P/L^2$ and $F=P/(B_1B_2)$ agree. Every specified sign
is strict, and $X$ is finite, so all of them persist for sufficiently
small positive $\varepsilon$.

For completeness, the following explicit inequalities suffice.
At every cube vertex require
\[
 \varepsilon\max(B_1,B_2)<L/2;
 \tag{3}
\]
at every point of $X$ where $(A_1+A_2)B_1B_2\ne0$, require
\[
 \varepsilon |(A_1+A_2)B_1B_2|<L|P|/2.
 \tag{4}
\]
If $l_i$ is a slope of $L$ and $b_{h,i}$ a slope of $B_h$, also
require $\varepsilon|b_{h,i}|<|l_i|/2$ whenever $b_{h,i}\ne0$.
All right-hand bounds are strictly positive, and there are finitely
many of them. Their minimum therefore admits a positive
$\varepsilon$. Equations (3) and the slope bounds prove the full-cube
positivity and strict orientation assertions. Equation (4) and (2)
prove the strict sign assertion directly.
\end{proof}

\paragraph{Biases and the Boolean threshold convention.}
A bias $c$ can first be absorbed by replacing $A_1$ with $A_1+cB_1$.
For a representation whose false inputs may have score zero, first
subtract a sufficiently small positive constant: finitely many true
inputs have a positive minimum score, so their signs are preserved,
while every false score becomes negative. The constant functions
cause no difficulty. Thus the theorem applies to the original strict
Boolean classification convention after this elementary finite repair.

\paragraph{Comparison with two-head coalescence.}
The earlier exact two-head coalescence uses one of the original
denominators as its limiting denominator. Formula (1) instead permits
an arbitrary prescribed limiting affine denominator and therefore
reverses its orientation when necessary, at a cost of one extra head.
The proof relies on the cleared numerator being a sum of two products.
For more heads the cleared numerator contains products of more affine
forms; replacing it by (1) one pair at a time does not preserve the
sign of the total score, because the positive multipliers differ.

\paragraph{Lean verification checkpoint.}
The module
\texttt{HeadComplexity/Results/TwoHeadOrientationTransfer.lean}
checks the exact identity as \texttt{three\_fraction\_identity}.
Its theorem \texttt{exists\_reverse\_three} proves simultaneous strict
sign preservation on the whole Boolean cube, positivity of every
output denominator on that cube, and strictly negative output slopes.
The theorem \texttt{reverse\_ratio\_exists\_atom} realizes each output
ratio by an original admissible atom with $\alpha<1$.
Finally, \texttt{exists\_reverse\_atoms\_of\_two} starts from the
unchanged proposition \texttt{fracComputable n 2 f}, strictifies its
threshold certificate, absorbs its bias, and produces three such
atoms with strict signs and zero external bias. These declarations
use no new axioms or unproved lemmas. The more general written
formulation for signs specified only on a finite subset, and for
the positive output orientation, is not claimed as Lean-checked by
this module. Neither version supplies an unbounded-head transfer
or the asymptotic improvement sought in this run.

\end{document}
